\originalchapter{8.05 作业7：动力学表象与相干态}
来源为8.05 Fall 2013 Assignment 7, 2013年10月24日发布, 11月1日15时截止, 七题共65分. 题面全部小问取自OCW原件, 解答为\textbf{AI独立解答}. 第1、2题$\lambda$是磁矩与自旋的比例系数, 第5题局部同名$\lambda$是展宽系数, 与其他章的无量纲扰动参数不同.
\originalsection{第1题：振荡磁场中的自旋（10分）}
\begin{exercise}
$\boldsymbol B=B_0\cos\omega t\,\hat z$, $H=-\lambda B_0\cos\omega t\,S_z$, 初态为$S_x=\hbar/2$. (a)求演化算符, 注意各时刻$H$对易. (b)求态及方向角. (c)求$S_x=-\hbar/2$概率. (d)求允许完全翻转的最大频率.
\end{exercise}
\begin{solution}
令$\chi(t)=\lambda B_0\sin\omega t/\omega$（$\omega=0$按连续极限）. 因为所有$H$都是$S_z$的实倍数,$U=\exp[\ii\chi S_z/\hbar]$. 从$\ket{x;+}=(\ket++\ket-)/\sqrt2$演化为$(\e^{\ii\chi/2}\ket++\e^{-\ii\chi/2}\ket-)/\sqrt2$. 自旋方向$\theta=\pi/2$, $\phi=-\chi$模$2\pi$, 并非磁场振幅本身作为角速度积分后仍为正. 翻转概率为$\sin^2(\chi/2)$. 完全翻转要求$|\chi|$能达到$\pi$, 所以正频率最大为$\omega_{\max}=|\lambda B_0|/\pi$. 若磁矩或场为0, 不存在完全翻转; 可使多个周期内相位来回变化, 但不能累计突破这个振幅界.
\end{solution}
\originalsection{第2题：Heisenberg自旋算符（5分）}
\begin{exercise}
常场$H=-\lambda BS_z$, 求$S_{xH}(t),S_{yH}(t),S_{zH}(t)$.
\end{exercise}
\begin{solution}
$O_H=U^\dagger OU$, 或直接$\dot S_{xH}=\lambda BS_{yH}$,$\dot S_{yH}=-\lambda BS_{xH}$. 故$S_{xH}=S_x\cos\lambda Bt+S_y\sin\lambda Bt$,$S_{yH}=S_y\cos\lambda Bt-S_x\sin\lambda Bt$,$S_{zH}=S_z$. 算符时间变, Heisenberg态固定; 与Schrödinger表象得到相同实验均值.
\end{solution}
\originalsection{第3题：均值与Newton方程（10分）}
\begin{exercise}
(a)对$H=p^2/(2m)+V(x)$推Heisenberg方程、Ehrenfest定理及经典化条件. (b)自由初态均值$x_0,p_0$, 求$\langle x\rangle(t)$. (c)改$V=qE_0x\sin\omega t$, 求不同时刻$H$对易子并解释Ehrenfest仍成立. (d)求新位置均值.
\end{exercise}
\begin{solution}
由$\dot O_H=U^\dagger(\ii[H,O]/\hbar)U$得$\dot x_H=p_H/m$, $\dot p_H=-V'(x_H)$; 对固定初态取平均得$\ddot{\langle x\rangle}=-\langle V'(x)\rangle/m$. 若$V'$是仿射函数, 精确等于$-V'(\langle x\rangle)/m$; 窄波包且力变化缓慢时可近似. 自由时$\langle x\rangle=x_0+p_0t/m$.

含时势给$[H(t_1),H(t_2)]=\ii\hbar qE_0[\sin\omega t_1-\sin\omega t_2]p/m$, 一般非零; 题面的“不同时刻非零”有$\sin\omega t_1=\sin\omega t_2$等特例, 应按一般不对易理解. 演化方程仍是$\ii\hbar\dot U=H(t)U$, 不需假定普通无时间排序的指数, 因而均值定理仍成立. 两次积分得到
\[
\langle p\rangle=p_0+\frac{qE_0}{\omega}(\cos\omega t-1),\qquad
\langle x\rangle=x_0+\frac{p_0t}m+\frac{qE_0}{m\omega^2}(\sin\omega t-\omega t).
\]
这里题给势的力为$-qE_0\sin\omega t$, 不能改成相反号. $\omega\to0$时力也趋0, 表达式连续还原自由均值.
\end{solution}
\originalsection{第4题：三维维里定理（10分）}
\begin{exercise}
中心势$H=\boldsymbol p^2/(2m)+V(r)$, $\Omega=\boldsymbol r\cdot\boldsymbol p$. (a)求$\dot\Omega_H$. (b)解释定态下时间无显式依赖算符的均值恒定. (c)对$V=c/r^k$证$\langle T\rangle=-(k/2)\langle V\rangle$.
\end{exercise}
\begin{solution}
逐方向应用对易规则得$\dot\Omega_H=\boldsymbol p_H^2/m-\boldsymbol r_H\cdot\boldsymbol\nabla V(r_H)=2T_H-r_HV'(r_H)$. 定态的Schrödinger波函数只乘相位, 所以任意固定观测量均值不变; 表象等价给$\mathrm d\langle\Omega_H\rangle/\mathrm dt=0$. 从而$2\langle T\rangle=\langle rV'\rangle=-k\langle V\rangle$. 原题$\Omega$不是Hermitian但定态对易平均仍可用; 改对称的$(\boldsymbol r\cdot\boldsymbol p+\boldsymbol p\cdot\boldsymbol r)/2$更适合作为观测量. 对奇异势须检查正规束缚态与原点边界条件, 不是任意$k,c$都自动存在这种态.
\end{solution}
\originalsection{第5题：恒力下的展宽（10分）}
\begin{exercise}
$H=p^2/(2m)+gx$. (a)积分Heisenberg方程给$x_H,p_H$. (b)初态$\psi=N\e^{-x^2/(2\Delta^2)}$,$\Delta>0$, 求位置均值. (c)求位置方差二次增长系数, 说明其对$g$的依赖.
\end{exercise}
\begin{solution}
$p_H=p-gt$, $x_H=x+pt/m-gt^2/(2m)$, 所以$\ddot x_H=-g/m$. 初态两均值均0, $\langle x_H\rangle=-gt^2/(2m)$, 与经典同初值一致. 方差中确定的$gt^2/(2m)$消去, 得
\[
(\Delta x(t))^2=(\Delta x(0))^2+\frac{t}{m}\langle xp+px\rangle_0+\frac{t^2}{m^2}(\Delta p(0))^2
=\frac{\Delta^2}{2}+\frac{\hbar^2t^2}{2m^2\Delta^2}.
\]
实偶Gaussian对称相关矩为0, 所以$\lambda=\hbar^2/(2m^2\Delta^2)$. 均匀力只加整体动量偏移、改变中心轨迹, 不改该展宽系数; 对空间不均匀力则不成立.
\end{solution}
\originalsection{第6题：位移谐振子（10分）}
\begin{exercise}
$H=p^2/(2m)+m\omega^2x^2/2-Fx$, $F>0$. (a)求基态能量和位置均值. (b)把新基态写作$N\e^{\alpha a^\dagger}\ket0$, 求实$\alpha$.
\end{exercise}
\begin{solution}
完成平方, 令$x_c=F/(m\omega^2)$,$y=x-x_c$, 得$H=p^2/(2m)+m\omega^2y^2/2-F^2/(2m\omega^2)$. 故$E_g=\hbar\omega/2-F^2/(2m\omega^2)$, $\langle x\rangle=x_c$. 新湮灭算符$a_y=a-\sqrt{m\omega/(2\hbar)}x_c=a-\alpha$. 要求$a_y\ket{0'}=0$给$\alpha=F/\sqrt{2m\hbar\omega^3}$. 用$a\e^{\alpha a^\dagger}\ket0=\alpha\e^{\alpha a^\dagger}\ket0$检验. 若归一且相位取实正,$N=\e^{-\alpha^2/2}$.
\end{solution}
\originalsection{第7题：相干态位置波函数（10分）}
\begin{exercise}
$\ket\alpha=\exp(\alpha a^\dagger-\alpha^*a)\ket0$,$\alpha=x_0/(\sqrt2d)+\ii p_0d/(\sqrt2\hbar)$,$d=\sqrt{\hbar/(m\omega)}$. 求显式归一化波函数.
\end{exercise}
\begin{solution}
指数代入$x,p$给$\alpha a^\dagger-\alpha^*a=\ii(p_0x-x_0p)/\hbar$. 因其两项对易子为常数, BCH给$D=\e^{\ii p_0x/\hbar}\e^{-\ii x_0p/\hbar}\e^{-\ii p_0x_0/(2\hbar)}$. 平移算符将$\psi_0(x)$变为$\psi_0(x-x_0)$, 故
\[
\psi_\alpha(x)=\frac1{(\pi d^2)^{1/4}}\exp\!\left[-\frac{(x-x_0)^2}{2d^2}+\frac{\ii p_0}{\hbar}\left(x-\frac{x_0}{2}\right)\right].
\]
模方是宽$d/\sqrt2$的归一Gaussian, 平均位置$x_0$、动量$p_0$, 相位常数$-p_0x_0/2$来自题给位移算符约定, 可用BCH核对而非随意省略.
\end{solution}
